\section{Limitations}
\label{sec:limitations}

% Scope bounds on conclusions. Observed evidence in Results remains valid within these bounds.
% Evidence: canonical protocol (seed null); paper/analysis.md; paper/conclusion.md; Method/Setup/Discussion.

The observations in Sections~\ref{sec:results}--\ref{sec:discussion} remain valid as
single-run records under the frozen protocol.
The limitations below restrict the \emph{scope of conclusions}, not the existence of those
recorded outcomes.

\subsection{Single-seed Evaluation}
\label{sec:lim-seed}

All primary configurations are single training trajectories with unset seed
($\texttt{seed}{=}\texttt{null}$ in the canonical protocol; Section~\ref{sec:setup}).
We do not report mean$\pm$std over seeds, and we do not perform formal significance tests.
Consequently, small deltas must not be read as statistically validated effects.
In particular, the Ablation~A vs.\ ResNet-20-A gap of $\DeltaAblationAMinusResNet$\,pp
(Section~\ref{sec:results-ablation}) is an observed single-run difference only.

\subsection{Frozen Training Recipe}
\label{sec:lim-recipe}

PreAct-ResNet-20-A and both ablations reuse the frozen post-activation baseline recipe
(learning rate, milestones, epochs, augmentation, Option~A, depth).
No PreAct-specific learning-rate or schedule retuning is present in the evidence bundle.
The negative PreAct comparison therefore answers whether a \emph{drop-in} pre-activation
replacement improves ResNet-20-A \textbf{under the frozen recipe}
($\ErrPreAct\%$ vs.\ $\ErrResNet\%$; $\DeltaPreActMinusResNet$\,pp).
It does \emph{not} answer whether a sufficiently retuned PreAct model would outperform
ResNet-20-A.

\subsection{Limited Architecture and Dataset Scope}
\label{sec:lim-scope}

All reported runs use CIFAR-10, depth~20, and Option~A residual shortcuts within the
present residual implementation (Sections~\ref{sec:method}--\ref{sec:setup}).
We do not claim ImageNet results, multi-depth Table~6 coverage, or transfer to other
datasets or network families.
Reproduction status for the residual baseline remains \ReproStatus at the depth-20 subset
tolerance $\ReproTolerancePP$\,pp (canonical reference), not a full-suite reproduction.

\subsection{Ablation~B Failure}
\label{sec:lim-ablation-b}

Ablation~B failed with numerical instability under the current recipe
(status \AblationBStatus) and therefore has no comparable best test error
(Section~\ref{sec:results-ablation}).
This limits accuracy comparisons involving Ablation~B; it does not license a universal claim
that the second BN is necessary in all settings.
A controlled rescue experiment (for example, lower learning rate) is not part of the present
SoT
\texttt{[NEEDS\_VERIFICATION]} / future work.

\subsection{Reproducibility and Environment Scope}
\label{sec:lim-env}

The evidence is a registered set of single-run artifacts (manifests, results, and the
canonical snapshot), not a multi-seed replication study.
Hardware and software notes for the frozen baseline protocol include an
NVIDIA GeForce RTX~3050 Laptop GPU (4\,GB) and PyTorch $2.0.1{+}\mathrm{cu117}$ under
conda \texttt{deeplearn} (Section~\ref{sec:setup}).
Exact Phase~3/5 driver/CUDA string identity is not exhaustively audited in the present paper
package
\texttt{[NEEDS\_VERIFICATION]};
we do not invent additional seed, driver, or CUDA details beyond the recorded sources.
